\scanpage{6}
- \mspage{3} -
\diagram{\begin{tikzcd}
p_*\Omega^{\mathrm{reg}}_{E/S}\arrow[r,"\sim"]&\unsure{o^*|e^*}\Omega^1_{E/S}=\omega
\end{tikzcd},}
et que pour la flèche adjointe à la flèche inverse
\diagram{\begin{tikzcd}
p^*\omega\arrow[r,"\sim"]&\Omega^{\mathrm{reg}}_{E/S}
\end{tikzcd}.}
Les sections de $\Omega^{\mathrm{reg}}_{E/S}$ ainsi associées aux sections de $\omega$ s'appellent les différentielles invariantes.

Soit $\pi$ une section inversible de $\omega$ ; il existe, localement pour la topologie de Zariski, une base $(y,x,1)$ de $p_*\mathcal O(3e)$ telle que l'image de $y$ dans $\omega^{\otimes-3}$ soit $-\pi^{\otimes-3}$, que $x$ appartienne à $p_*\mathcal O(2e)$, et que son image dans $\omega^{\otimes-2}$ soit $\pi^{\inserted{\otimes}-2}$. Si $\pi'$ est une autre section inversible de $\omega$, les bases correspondantes s'écrivent toutes, sous forme implicite (pour $r,s,t$ convenables)
\margin{G?}
\[
(1.2)\qquad\left\{\begin{aligned}
x&=u^2x'+r\\
y&=u^3y'+su^2x'+t\\
\pi'&=u\pi
\end{aligned}\right.\qquad .
\]
Le faisceau $\mathcal O(3e)$ est très ample et $(x,y,1)$ sont les coordonnées projectives d'un plongement de $E$ dans $P_2/S$.

La courbe image vérifiera une équation
\[
(1.3)\qquad y^2+a_1xy+a_3y=x^3+a_2x^2+a_4x+a_6.
\]
Pour vérifier l'existence et l'unicité de cette équation, on
\margin{5}
